\documentclass[10pt,a4paper]{amsart}
\usepackage[margin=19mm]{geometry}
\usepackage{amsmath,amssymb,mathtools}
\usepackage[scr=rsfs,cal=euler]{mathalfa}
\usepackage{xfrac}
\usepackage[unicode,hidelinks,bookmarks=false]{hyperref}
\newcommand{\ie}{i.e.\ }
\newcommand{\cf}{cf.\ }
\newcommand{\resp}{respectively\ }
\newcommand{\Subsection}{\subsection}
\providecommand{\nobreakdash}{\nobreak\hbox{-}}
\setlength{\emergencystretch}{2em}
\multlinegap0pt
\usepackage{bm}
\usepackage{bbm}
\usepackage{dsfont}
\usepackage{xspace}
\usepackage{cancel}
\usepackage{mathtools}
\usepackage{amsthm}
\makeatletter
\@ifundefined{lem}{\newtheorem{lem}{Lemma}}{}
\makeatother
\title[Arc weighted acyclic orientations and variations of degenera]{Arc weighted acyclic orientations and variations of degeneracy of graphs}
\author{Huan Zhou, Jialu Zhu, Xuding Zhu}
\date{Source: arXiv:2308.15853v4 (CC BY 4.0)}
\begin{document}
\maketitle

\noindent\emph{Source and licence.} Arc weighted acyclic orientations and variations of degeneracy of graphs. Version arXiv:2308.15853v4 (CC BY 4.0), distributed under the Creative Commons Attribution 4.0 International licence, as recorded in the arXiv licence metadata for this exact version and shown on the abstract page https://arxiv.org/abs/2308.15853v4. Full rights notice: licenses/arxiv-2308.15853-CC-BY-4.0.txt.\par\smallskip

\noindent\textit{Editorial scope.} Counted unit: the Lem whose source label is \texttt{lem-1} in the article (source0.tex lines 442--445, source label \texttt{lem-1}), delivered together with the article's own complete original proof of it (source0.tex lines 446--511, one \texttt{proof} environment of 66 lines). No mathematics is written, repaired or re-derived: statement and proof are the article's own text line for line. Required context retained uncounted: none. Every definition, display and label of those lines is retained unchanged. Omitted from this package and not redistributed: 1--441, 512--1041. Nothing inside a retained excerpt is omitted or altered: every retained body is delivered exactly as the article's lines are written, so no omission receipt of any kind accompanies this package. Citations: the retained excerpts carry no citation command at all, so no bibliography entry of the article is transcribed and no reference is redistributed. Reference adaptation: none was needed and none was made; every reference inside the delivered unit resolves inside it, so no unresolved reference or citation is produced. Printed slips kept literal: no printed word, symbol, hypothesis or number of the counted statement or of its proof was altered. Layout and class: the article's own class is \texttt{article} with packages including mathrsfs, amsmath, amsthm, amssymb, bm, mathtools, thm-restate, ulem, cancel, soul, geometry, fontenc, enumitem, natbib; the wrapper is the lane's pinned \texttt{amsart} editorial wrapper at 10pt on A4, which restates the theorem-like environments the retained text needs and adds the article's own macro definitions that the retained text needs (none), with extra wrapper packages (bm, bbm, dsfont, xspace, cancel, mathtools). The delivered numbering is the unit's own. Assets: no retained span contains a figure environment, a graphics-inclusion command, a PDF-page inclusion, a table or a TikZ picture; no image file is copied into this package, the package declares no asset dependency and no image file is redistributed with it. Licence: version arXiv:2308.15853v4 of this article is distributed under the Creative Commons Attribution 4.0 International licence, as recorded in the arXiv licence metadata for this exact version and shown on the abstract page https://arxiv.org/abs/2308.15853v4. A full rights notice is delivered at licenses/arxiv-2308.15853-CC-BY-4.0.txt. Characters: every retained excerpt is pure ASCII, so the delivered engine, XeLaTeX with its default Latin Modern text font, reports no missing character.
		\begin{lem}
			\label{lem-1}
			Let $G$ be a graph, $f,g \in \mathbb{N}^G$ and $0 \ne g \le f+1$. If $G$ is arc-weighted  $f$-degenerate, then  there is a non-empty subset $X$ of $V(G)$ such that $G[X]$ is arc-weighted  $(g-1)$-degenerate and $G-X$ is arc-weighted $(f-g)$-degenerate. 
		\end{lem}

		\begin{proof}
	The proof is by induction on $\sum_{v \in V(G)}f(v)$. 

 % \xztodo{Originally, Let $0 \ne g \le f$. If $G$ is AW strict $f$-degenerate, then $G[X]$ is strict AW $g$-degenerate, and $G-X$ is AW strict $(f-g)$-degenerate.
  %This means that if $G$ is AW $(f-1)$-degenerate, then $G[X]$ is AW $(g-1)$-degenerate, and $G-X$ is AW $(f-g-1)$-degenerate. Let $f'=f-1$ and $g'=g$. This becomes if $G$ is AW $f'$-degenerate, then $G[X]$ is $(g'-1)$-degenerate, and $G-X$ is AW $(f'-g')$-degenerate. If we choose to let $g'= g-1$, then $g'(v)=-1$ if $g(v)=0$. (We do allow $g(v)=0$, which means $v$ is not chosen in this move, and no colour is given to $v$). So $g'=g$ still represents the number of colours. Hence degeneracy becomes $(g-1)$-degenerate.}
 
 Assume that $G$ is arc-weighted $f$-degenerate. 
 If $E(G)=\emptyset$, then let $X=\{v: g(v)  \ge 1\}$, and the conclusion holds.  
 
  Assume $|E(G)| \ge 1$, and $(D,w)$ is a  minimal arc-weighted   acyclic orientation of $G$. If there is a vertex $x$ for which  
$f(x) > d_{(D,w)}^+(x)$, 
then let $f'=f-1_{\{x\}}$ and $g' = \max\{0, g- 1_{\{x\}}\}$. Then $g' \le f'+1$ and $G$ is arc-weighted  $f'$-degenerate. If $g'=0$, then $g=1_{\{x\}}$.  
  Let $X=\{x\}$. From the definition it follows that $G[X]$ is arc-weighted $(g-1)$-degenerate and $G-X$ is arc-weighted $(f-g)$-degenerate. 

   Otherwise, assume $g' \neq 0$. By the induction hypothesis,  there is a non-empty  subset $X$ of $V(G)$, such that $G[X]$ is arc-weighted   $(g'-1)$-degenerate, and $G-X$ is arc-weighted   $(f'-g')$-degenerate. Then $G[X]$ is arc-weighted  $(g-1)$-degenerate, and $G-X$ is arc-weighted $(f-g)$-degenerate.

Assume $f(x) = d_{(D,w)}^+(x)$ for each vertex $x$. 
By definition, $(D,w)$ has a dominating arc $e=(x,y)$.  As $(D,w)$ is a  minimal arc-weighted   acyclic orientation,  $w(e) = d_{(D,w)}^+(y)+1 =f(y)+1$.
Then $G'=G-e$ is  arc-weighted   $f'$-degenerate, where $f'=f-(f(y)+1)1_{\{x\}}$.
 

    Let $g' \in \mathbb{N}^{G'}$ be defined as $g'(v)=g(v)$, except that $$g'(x)=\min\{f'(x)+1, \max\{0,g(x)-g(y)\}\}.$$ 
   If $g'=0$, then  $g(x) \le g(y)$ and $g(v)=0$ for $v \ne x$. But $g(x) \le g(y) = 0$ implies that $g=0$, a contradiction.
   Thus $g' \ne 0$. By the induction hypothesis,  there is a non-empty   subset $X$ of $V(G')$, so that 
$G'[X]$ is arc-weighted  $(g'-1)$-degenerate, and $G'-X$ is arc-weighted   $(f'-g')$-degenerate. 

\bigskip 
\noindent
{\bf Case 1}  $x \in X$. 



As $G-X=G'-X$ is arc-weighted   $(f'-g')$-degenerate, and $(f-g)(v)= (f'-g')(v)$ for $v \in V(G-X)$,   we conclude that  $G-X $ is arc-weighted  $(f-g)$-degenerate. 
It remains to show that $G[X]$ is arc-weighted $(g-1)$-degenerate.

If $y \notin X$, then
$G[X]=G'[X]$ is arc-weighted  $(g'-1)$-degenerate, and hence arc-weighted $(g-1)$-degenerate, since $g'(v)\le g(v)$ for every $v\in X$. 


Assume that $y \in X$.
 As $G'[X]$ is arc-weighted   $(g'-1)$-degenerate, and $x \in X$, we have   $g'(x) \ge 1$.  Hence 
$g'(x) = \min\{f'(x)+1, g(x)-g(y)\} \le g(x)-g(y)$.  Let $(D',w')$ be  an arc-weighted   acyclic orientation  of $G'[X]$ with $d_{(D',w')}^+  \le g'-1$. 
 Let $(D'',w'')$ be obtained from $(D',w')$ by adding the arc $e=(x,y)$ of weight $w''(e)= g(y)=g'(y)$. Then $$w''(e)   = g'(y)   \ge d_{(D',w')}^+(y)+1 = d_{(D'',w'')}^+(y)+1,$$
and $$d_{(D'',w'')}^+(x) = d_{(D',w')}^+(x)+w''(e) < g'(x)+g(y) \le  g(x).$$
Hence, $e$ is a dominating arc in $(D'',w'')$ and $(D''-e,w'') = (D',w')$. Therefore $(D'',w'')$ is an arc-weighted acyclic orientation of $G[X]$ with 
$ d_{(D'',w'')}^+  \le  g-1$. Hence $G[X]$ is arc-weighted   $(g-1)$-degenerate. 



\bigskip 
\noindent
{\bf Case 2}   $x \notin X$. 

Then  $G[X] = G'[X]$ is arc-weighted  $(g'-1)$-degenerate,  and $g'(v)=g(v)$ for all $v \in X$. Hence $G[X]$ is arc-weighted   $(g-1)$-degenerate. 
It remains to show that $G-X$ is arc-weighted $(f-g)$-degenerate.

As $G'-X$ is arc-weighted $(f'-g')$-degenerate,  we know that $(f'-g')(x) \ge 0$. So $g'(x) \ne f'(x)+1$ and hence $$ g'(x)= \max\{0, g(x)-g(y)\} \ge g(x)-g(y).$$ 
This implies that $(f'-g')(x) \le (f(x)-(f(y)+1)) - (g(x)-g(y)) \le (f-g)(x)$   (note that $f(y)+1 \ge g(y)$). 
If   $y \in X$,   then   $G-X = G'-X$ is arc-weighted  $(f-g)$-degenerate. 
  Therefore, we assume that $y \notin X$. 

  Let $(D',w')$ be an arc-weighted  acyclic orientation of $G'-X$ with $d_{(D',w')}^+   \le (f'-g')$.  
  Let $(D'',w'')$ be obtained from $(D',w')$ by adding the arc $e=(x,y)$ of weight $w''(e)=f'(y)-g'(y)+1=f(y)-g(y)+1$. Thus, $(D'',w'')$ is an arc-weighted   acyclic orientation of $G-X$. As $d_{(D'',w'')}^+(x)   \le (f'-g')(x) + w''(e) \le (f-g)(x)$, we have  $d_{(D'',w'')}^+  \le f-g$, and hence $G-X$ is arc-weighted   $(f-g)$-degenerate.
 
   This completes the proof of Lemma \ref{lem-1}.
		\end{proof}	

\end{document}
